\chapter{Introduction}
\label{ch:introduction}

Online assessment in higher education has shifted from \emph{convenient} to \emph{essential}. A modern computer-science classroom expects students to write code in real time, answer multiple-choice questions about its complexity, and explain their reasoning in prose --- all in one sitting, with sandboxed execution, automatic grading where it is safe, and a workflow that scales from a single classroom to a department. Existing learning-management systems have a long history of bolting these capabilities onto a quiz module; existing online judges have a long history of treating classroom organisation as somebody else's problem. The result, as \Cref{ch:literature} discusses, is a workflow where the teacher integrates two or three tools manually --- and the student feels the seams.

\section{Problem statement}

Three concrete pain points motivated this project.

\paragraph{One exam, three formats.} A computer-science instructor frequently wants a single timed exam that mixes a coding problem (to assess implementation), an MCQ section (to assess conceptual understanding), and a written question (to assess explanation). Today, those three lives in different tools or different modules, with separate timers, separate gradebooks, and separate notification streams. Students switch tools mid-exam, teachers reconcile scores across exports, and small operational details (a network blip during the coding section) become large operational pains.

\paragraph{Auto-grading that does not embarrass itself.} Auto-grading by I/O diff is a thirty-year-old technique~\cite{douce2005autoassess}, but the standard implementations still trip on whitespace differences and trailing newlines often enough that a working system has to handle the embarrassment by hand. A system that wants its grading to be trusted has to normalise output without compromising the detection of genuinely wrong answers, and has to surface compilation errors and runtime errors as separate verdicts so a student can act on the diagnostic.

\paragraph{Self-hosting without operational pain.} A large institution can buy access to a commercial platform; a department or a single instructor cannot. The available open-source pieces (Moodle, Judge0, Canvas Open Source) compose, but the composition is itself a project. A self-hosted alternative that is "one Go binary plus PostgreSQL plus Judge0" is the gap APEX targets.

\section{Aim and objectives}

The aim of this graduation project is to design and build APEX: a self-hostable web portal that integrates coding, MCQ, and written exam delivery, with sandboxed auto-grading for code, automatic grading for MCQ, manual grading for written answers, and a single classroom workflow that ties everything together.

The aim is decomposed into six objectives:

\begin{enumerate}
  \item[\textbf{O1}] Deliver mixed-format exams (coding, MCQ, written) in a single timed attempt, with autosave, resume, and a single submit button.
  \item[\textbf{O2}] Auto-grade coding submissions through a sandboxed Judge0~\cite{judge0,isolate} integration with normalised output diffing.
  \item[\textbf{O3}] Manually grade written submissions through a queue, with score, status, and textual feedback that re-aggregate the attempt's score.
  \item[\textbf{O4}] Authenticate users with email + password and Google OAuth, defend against the four common classes of weakness in classroom portals (insecure direct object reference, stored cross-site scripting, weak fall-back secrets, HTML injection in email), and re-check ownership at the handler level.
  \item[\textbf{O5}] Deploy as a single Go binary plus PostgreSQL, with optional Judge0 and SMTP, configurable through environment variables, with idempotent schema migrations.
  \item[\textbf{O6}] Provide a usable interface for both teacher and student roles, evaluated against the System Usability Scale~\cite{brooke1996sus} with a small user study.
\end{enumerate}

Each objective maps directly onto the requirements in \Cref{ch:requirements} and onto the implementation in \Cref{ch:implementation}.

\section{Approach in brief}

APEX is built as a three-tier web application: a React + Vite + TypeScript single-page application in the browser, a Go (Gin + GORM) backend exposing a typed REST API, and PostgreSQL as the durable store. Judge0 is a fourth-party sandbox; SMTP and Google Identity Services are optional. \Cref{ch:design} unpacks this in detail; \Cref{fig:F4.1} draws the deployment context.

The implementation totals approximately 10,000 lines of Go on the backend and 12,000 lines of TypeScript on the frontend. The schema has fourteen first-class entities, documented in \Cref{fig:F4.4}. A 60-second-tick reminder scheduler dispatches in-app reminders and exam-start emails; a security review (\Cref{ch:security}) hardened four classes of weakness before the system left development.

\Cref{fig:F1.2} positions APEX against four representative comparators: Moodle (with the CodeRunner plugin), Canvas LMS, HackerRank for Schools, and CodeRunner standalone. The complete comparison lives in \Cref{ch:literature}.

\begin{figure}[H]
  \centering
  \includegraphics[width=0.85\linewidth]{F1.2_feature-radar.pdf}
  \caption{Feature-coverage radar (subjective; 0--5, higher is better). APEX is positioned to cover both the LMS-style classroom features and the judge-style coding features at "good enough", at the cost of being newer and less battle-tested than any single comparator.}
  \label{fig:F1.2}
\end{figure}

\section{Scope and non-scope}

\paragraph{In scope.} Identity (signup, login, OAuth, password reset, account deletion); class management (creation, invite codes, members, announcements); exam authoring (mixed formats, scheduling, drafts, assignment, preview); exam taking (autosave, resume, sample-test runs, submission); grading (auto for coding and MCQ, manual queue for written); reminders (in-app and email); results (per-student, per-question, per-test detail); security (the patches in \Cref{ch:security}); deployment (single binary, idempotent migrations); usability evaluation (\Cref{ch:results}).

\paragraph{Out of scope.} Multi-tenant deployment; cross-region failover; commercial proctoring (webcam, screen capture); plagiarism detection (MOSS~\cite{moss}, JPlag~\cite{prechelt2002jplag} bridges are on the roadmap, not in scope for this iteration); integration with institutional SIS or LTI; payment, billing, and licence management.

\section{Contributions}

This book contributes:

\begin{itemize}
  \item A working, self-hostable portal that integrates three exam formats under a single attempt, a single grader pipeline, and a single notification stream.
  \item A focused security review and a published patch set (\Cref{ch:security}, commit \texttt{2eaac88}) covering IDOR, stored XSS, weak fall-back secrets, and email HTML injection, with the resulting code patterns documented and reused throughout.
  \item A documented design that maps each architectural decision to the specific code site that implements it, suitable as a reference for future maintainers.
  \item A reproducible build pipeline for the system itself (\texttt{start.sh}, Docker Compose Postgres, Vite proxy) and for this book (\texttt{book/Makefile}).
  \item A small but honest evaluation: latency and throughput benchmarks against a self-hosted Judge0, and a SUS-based usability study with both teacher and student participants.
\end{itemize}

\section{Book structure}

The remaining chapters are organised as follows.

\begin{itemize}
  \item \Cref{ch:literature} surveys online-assessment platforms, online judges, auto-grading approaches, and the security and usability literature.
  \item \Cref{ch:requirements} formalises the functional and non-functional requirements and maps them onto use cases.
  \item \Cref{ch:design} documents the architecture: deployment context, layered backend, frontend structure, data model, state machines, sequence diagrams, the API surface, and the migration pipeline.
  \item \Cref{ch:implementation} walks the code that fills in the design, file by file and line by line.
  \item \Cref{ch:security} documents the threat model, the four findings, and the patches.
  \item \Cref{ch:testing} documents the test strategy.
  \item \Cref{ch:deployment} documents the development and production deployment stories.
  \item \Cref{ch:user-manual} is a screenshot-led walkthrough of the running app.
  \item \Cref{ch:results} reports the evaluation results.
  \item \Cref{ch:conclusion} summarises lessons learned and the roadmap.
\end{itemize}

The appendices reproduce the API table (\Cref{appendix:api-reference}), the schema DDL (\Cref{appendix:schema-ddl}), the environment-variable table (\Cref{appendix:env-vars}), notable source listings (\Cref{appendix:source-listings}), the screenshot capture script (\Cref{appendix:screenshots-extra}), and the project Gantt (\Cref{appendix:gantt}).
